\documentclass[11pt]{article}

\usepackage[T1]{fontenc}
\usepackage{lmodern}
\usepackage{amsmath,amssymb,amsthm,mathtools}
\usepackage{array,booktabs}
\usepackage{float}
\usepackage{graphicx}
\usepackage{microtype}
\usepackage[margin=1in]{geometry}
\usepackage[hidelinks]{hyperref}

\newtheorem{theorem}{Theorem}
\newtheorem{lemma}{Lemma}
\newtheorem{proposition}{Proposition}
\newcommand{\Mhat}{\widehat M}
\newcommand{\R}{\mathbb{R}}
\newcommand{\C}{\mathbb{C}}
\newcommand{\Z}{\mathbb{Z}}
\newcommand{\Tr}{\operatorname{Tr}}
\newcommand{\vecop}{\operatorname{vec}}

\title{A 16-Vertex Counterexample to Path Minimality for\\
Laplacian Quantum-Walk Laziness}
\author{Xiaojian Zeng\thanks{Email: \texttt{xiaojianzeng@foxmail.com}}\\
PBC School of Finance, Tsinghua University, Beijing 100190, China}
\date{}

\begin{document}
\maketitle

\begin{abstract}
For a continuous-time quantum walk generated by the combinatorial Laplacian
$L_G$ of a finite simple graph, let $\Mhat(G)$ be the infinite-time Ces\`aro
average of the transition-probability matrix and set
$\tau(G)=\Tr\Mhat(G)$.  It was asked whether the path $P_n$ minimizes
$\tau$ among all graphs on $n$ vertices.  We give an exact counterexample.
On $\Z/16\Z$, start with the circulant having cyclic differences
$\{\pm2,\pm3,\pm4,\pm6\}$ and delete the edges $\{0,13\}$ and
$\{5,8\}$.  The resulting connected, unweighted, loopless simple graph $G$
has 62 edges and satisfies
\[
  \tau(G)=\frac{1673}{1192}
  <\frac{23}{16}=\tau(P_{16}),
  \qquad
  \tau(G)-\tau(P_{16})=-\frac{81}{2384}.
\]
The sign is certified in exact rational arithmetic by orthogonal projection
onto the commutant of $L_G$; a structurally independent full-eigenspace
calculation provides a numerical cross-check.  Thus the path is not a global
minimizer of long-time return over all unweighted graph topologies.
\end{abstract}

\noindent\textbf{Keywords:} continuous-time quantum walk; average mixing
matrix; graph Laplacian; algebraic graph theory; exact counterexample

\section{Introduction}

Let $G$ be a finite graph with combinatorial Laplacian $L_G=D_G-A_G$.
The continuous-time walk generated by $L_G$ has transition matrix
$U_G(t)=\exp(-itL_G)$.  Entrywise squared amplitudes generally do not
converge, but their infinite-time Ces\`aro average does.  This average mixing
matrix is a spectral graph invariant introduced and developed in the study
of continuous quantum walks~\cite{Godsil2013}.  Its trace measures the
average long-time tendency to return to the initial vertex; recent work calls
this trace the \emph{laziness} of the walk~\cite{MohanTamonXuZhan2026}.

Godsil, Guo and Sobchuk studied extremal diagonal behavior and tabulated
Laplacian trace minima through eight vertices~\cite{GodsilGuoSobchuk2023}.
Mohan, Tamon, Xu and Zhan subsequently developed further extremal results and,
after reporting data through nine vertices, asked whether $P_n$ is among the
minimizers over all $n$-vertex graphs
\cite[Section~8(ii)]{MohanTamonXuZhan2026}.  The motivating formula is
Godsil's exact average mixing matrix for the Laplacian walk on a path
\cite[Theorem~5.3]{Godsil2013}.

We answer that universal question in the negative with one connected graph
on 16 vertices.  The witness is specified both by a concise cyclic
construction and by a complete edge list.  Its strict deficit is an exact
rational number, so neither a finite observation time nor a floating-point
tolerance determines the conclusion.  The result concerns only this
counterexample: it does not classify the minimizers at order 16, identify the
smallest order of failure, or give an infinite family or replacement bound.

\section{Laplacian walk and long-time return}
\label{sec:model}

Throughout, $G=(V,E)$ is finite, undirected, unweighted, loopless and simple,
with $|V|=n$.  Disconnected graphs are allowed in the universal question,
although the witness below is connected.  On the single-excitation Hilbert
space $\C^V$, take
\[
  H=L_G=D_G-A_G,\qquad U_G(t)=e^{-itL_G},\qquad \hbar=1,
\]
with unit edge coupling.  There is no additional potential or continuum
boundary.  All quantities below are exact finite-dimensional quantities; no
perturbation series, thermodynamic limit or time discretization is used.

Let $\Lambda(G)$ be the set of distinct eigenvalues of $L_G$, and let
$E_\lambda$ denote the orthogonal projector onto the full
$\lambda$-eigenspace.  The spectral decomposition gives
\[
  U_G(t)=\sum_{\lambda\in\Lambda(G)}e^{-it\lambda}E_\lambda.
\]
Time averaging removes terms with unequal frequencies, hence
\begin{equation}
  \Mhat(G)_{uv}
  =\lim_{T\to\infty}\frac1T\int_0^T|U_G(t)_{uv}|^2\,dt
  =\sum_{\lambda\in\Lambda(G)}(E_\lambda)_{uv}^{2}.
  \label{eq:average-mixing}
\end{equation}
The full eigenspace projector is essential at a degenerate eigenvalue.  We
define
\begin{equation}
  \tau(G)=\Tr\Mhat(G)
  =\sum_{v\in V}\sum_{\lambda\in\Lambda(G)}(E_\lambda)_{vv}^{2}.
  \label{eq:tau}
\end{equation}
Consequently, $\tau(G)/n$ is the long-time return probability averaged over a
uniformly selected initial vertex.  Sources that use $e^{+itL_G}$ give the
same observable: because $L_G$ is real symmetric, changing the sign of time
complex-conjugates $U_G(t)$ and leaves all probabilities unchanged.

For $P_n$, let $J$ be the all-ones matrix and let $R$ reverse the vertex
order.  The published formula is
\begin{equation}
  \Mhat(P_n)=\frac{(n-1)J+(n/2)(I+R)}{n^2}.
  \label{eq:path-matrix}
\end{equation}
When $n$ is even, $R$ has no fixed vertex, and tracing
\eqref{eq:path-matrix} gives
\begin{equation}
  \tau(P_n)=\frac32-\frac1n,
  \qquad \tau(P_{16})=\frac{23}{16}.
  \label{eq:path-trace}
\end{equation}

\section{The 16-vertex witness}
\label{sec:witness}

Write the vertices as $\Z/16\Z$.  Begin with the undirected circulant
\begin{equation}
  E_0=\bigl\{\{x,x+s\}:x\in\Z/16\Z,\ s\in\{2,3,4,6\}\bigr\},
  \label{eq:circulant}
\end{equation}
where addition is modulo 16, and define
\begin{equation}
  E(G)=E_0\setminus\bigl\{\{0,13\},\{5,8\}\bigr\}.
  \label{eq:deleted}
\end{equation}
The two deleted step-3 edges are exchanged by the half-turn
$x\mapsto x+8$; their endpoints are not individually antipodal.  The base
circulant is 8-regular with 64 edges.  The deleted edges are disjoint, so $G$
has 62 edges and degree sequence $(8^{12},7^4)$, with degree-seven vertices
$0,5,8,13$.  It contains the spanning path
\[
  0,2,4,6,8,10,12,14,1,3,5,7,9,11,13,15,
\]
which proves connectivity.  The construction uses only unit, undirected
edges and introduces neither loops nor multiple edges.

Figure~\ref{fig:witness} displays the graph construction and the exact return
comparison.  For complete independent transcription, Table~\ref{tab:edges}
lists every edge once: the row indexed by $u$ gives precisely the neighbors
$v>u$.

\begin{figure}[H]
  \centering
  \includegraphics[width=0.96\linewidth]{figures/witness_graph.pdf}
  \caption{(a) The witness graph, with the two deleted step-3 edges shown as
  red dashes and the four degree-seven vertices in orange.  Present edges are
  colored by cyclic step.  (b) Uniform-start long-time return probabilities,
  plotted from zero and labeled by their exact values.}
  \label{fig:witness}
\end{figure}

\begin{table}[H]
\centering
\small
\caption{Complete zero-based edge list for $G$, grouped by the smaller
endpoint.  Empty terminal rows for 14 and 15 are omitted.}
\label{tab:edges}
\begin{tabular}{@{}c>{\raggedright\arraybackslash}p{0.35\linewidth}
                  c>{\raggedright\arraybackslash}p{0.31\linewidth}@{}}
\toprule
$u$ & all $v>u$ with $\{u,v\}\in E(G)$ & $u$ & all $v>u$ with $\{u,v\}\in E(G)$ \\
\midrule
0 & 2, 3, 4, 6, 10, 12, 14 & 7 & 9, 10, 11, 13 \\
1 & 3, 4, 5, 7, 11, 13, 14, 15 & 8 & 10, 11, 12, 14 \\
2 & 4, 5, 6, 8, 12, 14, 15 & 9 & 11, 12, 13, 15 \\
3 & 5, 6, 7, 9, 13, 15 & 10 & 12, 13, 14 \\
4 & 6, 7, 8, 10, 14 & 11 & 13, 14, 15 \\
5 & 7, 9, 11, 15 & 12 & 14, 15 \\
6 & 8, 9, 10, 12 & 13 & 15 \\
\bottomrule
\end{tabular}
\end{table}

\section{Exact commutant certificate}
\label{sec:certificate}

We use the commutant representation of average mixing
\cite{CoutinhoGodsilGuoZhan2018}.  The following formulation makes the exact
calculation explicit.  Matrices are vectorized by stacking columns, and
$D_v=e_ve_v^{\mathsf T}$ and $d_v=\vecop(D_v)$.

\begin{lemma}
\label{lem:commutant}
Let $L$ be a real symmetric $n\times n$ matrix.  If the columns of $B$ form a
basis for the vectorized commutant
\[
  \mathcal C(L)=\{X\in\R^{n\times n}:LX=XL\},
\]
then
\begin{equation}
  Q=B(B^{\mathsf T}B)^{-1}B^{\mathsf T}
  \label{eq:projector}
\end{equation}
is its Frobenius-orthogonal projector and
\begin{equation}
  \Mhat(L)_{uv}=d_u^{\mathsf T}Qd_v,
  \qquad
  \tau(L)=\sum_{v=0}^{n-1}d_v^{\mathsf T}Qd_v.
  \label{eq:diagonal-restriction}
\end{equation}
\end{lemma}

\begin{proof}
If $L=\sum_\lambda\lambda E_\lambda$, the Frobenius projection onto
$\mathcal C(L)$ is $X\mapsto\sum_\lambda E_\lambda X E_\lambda$.
Consequently,
\[
  \langle D_u,\sum_\lambda E_\lambda D_vE_\lambda\rangle_F
  =\sum_\lambda(E_\lambda)_{uv}(E_\lambda)_{vu}
  =\Mhat(L)_{uv}.
\]
Equation~\eqref{eq:projector} is the standard matrix of the orthogonal
projection onto the column space of a full-rank $B$.  Restricting its rows and
columns to the diagonal matrix units proves
\eqref{eq:diagonal-restriction}.
\end{proof}

For the Laplacian of the edge list in Table~\ref{tab:edges}, exact integer
arithmetic gives the characteristic polynomial
\begin{align}
\chi_L(x)={}&x(x-10)(x-8)(x-6)(x^2-20x+98)(x^2-16x+62)\notag\\
&\times(x^4-32x^3+368x^2-1768x+2874)\notag\\
&\times(x^4-32x^3+380x^2-1984x+3842).
\label{eq:charpoly}
\end{align}
The exact check $\gcd(\chi_L,\chi_L')=1$ shows that the spectrum is simple.
Thus the minimal polynomial has degree 16 and the commutant has the basis
\begin{equation}
  B=\bigl[\vecop(I)\ \vecop(L)\ \cdots\ \vecop(L^{15})\bigr].
  \label{eq:power-basis}
\end{equation}
Equivalently, if
\[
  C=I\otimes L-L^{\mathsf T}\otimes I,
\]
then $CB=0$, the 16 columns of $B$ are independent, and they span
$\ker C$.

The evaluation can be compressed to a $16\times16$ rational calculation.
Let
\begin{equation}
  \Gamma=B^{\mathsf T}B,
  \qquad
  R_{v,k}=(L^k)_{vv},\quad 0\le v,k\le15.
  \label{eq:gram-moments}
\end{equation}
Selecting the diagonal rows of $B$ in Lemma~\ref{lem:commutant} gives
\begin{equation}
  \Mhat(G)=R\Gamma^{-1}R^{\mathsf T},
  \qquad
  \tau(G)=\Tr\!\left(R\Gamma^{-1}R^{\mathsf T}\right).
  \label{eq:exact-evaluation}
\end{equation}
Every entry of $L$, $R$ and $\Gamma$ is an integer.  The accompanying exact
script verifies $CB=0$ and $\Gamma\Gamma^{-1}=I$ symbolically before applying
\eqref{eq:exact-evaluation}; no eigenvalue approximation enters this route.

\begin{proposition}
For the graph $G$ of \eqref{eq:circulant}--\eqref{eq:deleted},
\[
  \tau(G)=\frac{1673}{1192}.
\]
\end{proposition}

\begin{proof}
Constructing the integer Laplacian from Table~\ref{tab:edges}, forming the
power basis \eqref{eq:power-basis}, and evaluating
\eqref{eq:exact-evaluation} in $\mathbb Q$ gives
\[
  \Tr(R\Gamma^{-1}R^{\mathsf T})=\frac{1673}{1192}.
\]
The input edge table, factorization \eqref{eq:charpoly}, exact matrix guards
and rational output are included in the machine-readable supplement, together
with the short reconstruction program.
\end{proof}

A separate calculation constructs the edges from the cyclic rule rather than
from Table~\ref{tab:edges}, diagonalizes the resulting symmetric Laplacian,
groups full eigenspace projectors at two gap tolerances, and checks residuals,
orthogonality, projector completeness and idempotence.  Both tolerances find
16 distinct spectral groups and give
$1.4035234899328854$, differing from $1673/1192$ by
$4.44\times10^{-16}$.  This floating-point route is a cross-check only; the
strict inequality below rests on the rational calculation.

\section{Disproof and physical interpretation}

Combining the proposition with \eqref{eq:path-trace} gives the exact identity
\begin{equation}
  \tau(G)-\tau(P_{16})
  =\frac{1673}{1192}-\frac{23}{16}
  =-\frac{81}{2384}<0.
  \label{eq:deficit}
\end{equation}

\begin{theorem}
The statement $\tau(X)\ge\tau(P_n)$ for every $n\ge2$ and every undirected,
unweighted, loopless simple $n$-vertex graph $X$ is false.
\end{theorem}

\begin{proof}
The graph $G$ in Section~\ref{sec:witness} belongs to the quantified class,
and \eqref{eq:deficit} violates the asserted inequality at $n=16$.
\end{proof}

For a uniformly chosen initial vertex, division by 16 gives
\[
  \frac{\tau(G)}{16}=\frac{1673}{19072},
  \qquad
  \frac{\tau(P_{16})}{16}=\frac{23}{256},
  \qquad
  \frac{\tau(G)-\tau(P_{16})}{16}=-\frac{81}{38144}.
\]
Thus this unweighted topology has a strictly smaller long-time return than
the path in the same finite, closed single-excitation Laplacian model.  The
claim is not that the propagator converges pointwise, nor that a finite-time
experiment has been performed; it is an exact statement about its Ces\`aro
average.

\section{Scope and outlook}

One strict admissible witness is enough to settle the universal
path-minimization question, but it leaves the revised extremal problem open.
In particular, the calculation does not establish the minimum at $n=16$, the
smallest order at which a counterexample exists, an equality classification,
an infinite counterexample family, or a new all-order lower bound.  The
available literature comparison is also bounded, especially because the
motivating question is recent; priority should therefore be rechecked before
submission.  A natural next mathematical task is to determine whether the
near-circulant mechanism in \eqref{eq:deleted} extends to a controlled family
or instead produces an isolated finite-order exception.

\section*{Data and code availability}

The local supplementary package contains the complete edge data, the exact
rational commutant calculation, the independent grouped-eigenspace
cross-check, the comparison script, the generated figure source and explicit
regeneration instructions.  The exact certificate requires only integer and
rational linear algebra.  No downloaded third-party paper is included in the
supplement.

\section*{AI-assistance disclosure}

Generative AI assisted literature triage, code generation, exact-calculation
workflows, verification and manuscript drafting.  The named human author must
independently check all claims, calculations, citations and wording, approve
the authorship metadata, and choose final disclosure language before any
submission.  This local draft does not represent completed human review.

\bibliographystyle{abbrv}
\bibliography{manuscript/references}

\end{document}
